% GENERATED FILE. Edit canonical lessons, then run npm run generate:books.
% canonical-source-hash: fnv1a64:307ab50ec2f384ff
% canonical-lessons: GE-C388-kredit, GE-C388-darlehen, GE-C388-schulden, GE-C388-zurueckzahlen, GE-C388-borgen

\chapter{Geld leihen}
\label{ch:ge-geld-leihen}
\hlchaptermodality{388}

\begin{chapteropening}
\textbf{By the end of this chapter:} \emph{I can talk about loans and debts, paying money back, and borrowing from friends.}

Complete the last lesson of chapter 388: I can talk about loans and debts, paying money back, and borrowing from friends.
\end{chapteropening}

\section[der Kredit]{der Kredit — the bank loan, the credit}

\label{lesson:GE-C388-kredit}



\begin{quote}
Before the new one: say the German for to debit, then the German for the amount.
\end{quote}

\subsection*{You'll want to know: der Kredit}
\textbf{der Kredit} — \textquotedblleft{}the bank loan, the credit\textquotedblright{}.

\begin{quote}
Für das Haus haben sie einen Kredit bei der Bank aufgenommen.
\end{quote}

\begin{quote}
\emph{They took out a bank loan for the house.}
\end{quote}

Plural: \emph{die Kredite}. \emph{Einen Kredit aufnehmen} — to take out a loan.

\begin{cousinweb}
\emph{Kredit} was borrowed from French \emph{crédit} (EtymDB), which goes back through Italian \emph{credito} to Latin \emph{crēditum} \textquotedblleft{}a loan\textquotedblright{}, from \emph{crēdere} \textquotedblleft{}to trust, to believe\textquotedblright{} (Lewis \& Short).
\end{cousinweb}

\begin{grammarlens}[title={What you've built}]
One more word for telling, explaining and describing at B1.
\end{grammarlens}

\subsection*{Your turn}
\begin{itemize}
\raggedright
  \item \emph{Say it:} \emph{der Kredit}
  \item \emph{Say it:} \emph{der Kredit}, once more
  \item \emph{Say it:} a sentence of your own with \emph{der Kredit}
  \item \emph{Recall:} say the German for to debit, then the German for the amount, then say \emph{der Kredit} again
\end{itemize}

\begin{tcolorbox}[breakable,colback=teal!4,colframe=teal!35!black,title={Before you move on}]
What does der Kredit mean? (\textquotedblleft{}The bank loan, the credit\textquotedblright{}.) Say it once more.
\end{tcolorbox}
\section[das Darlehen]{das Darlehen — the loan}

\label{lesson:GE-C388-darlehen}



\begin{quote}
Before the new one: say the German for the amount, then the German for the bank loan.
\end{quote}

\subsection*{You'll want to know: das Darlehen}
\textbf{das Darlehen} — \textquotedblleft{}the loan\textquotedblright{}.

\begin{quote}
Meine Eltern haben mir ein Darlehen für das Studium gegeben.
\end{quote}

\begin{quote}
\emph{My parents gave me a loan for my studies.}
\end{quote}

Plural: \emph{die Darlehen}. Formal and legal; in everyday speech \emph{der Kredit} is more common.

\begin{cousinweb}
\emph{Darlehen} joins \emph{dar} \textquotedblleft{}there, forth\textquotedblright{} to \emph{Lehen} \textquotedblleft{}a loan, a fief\textquotedblright{}, of the family of \emph{leihen} \textquotedblleft{}to lend\textquotedblright{}, Old High German \emph{līhan}, Proto-Germanic \emph{līhwaną} (EtymDB).
\end{cousinweb}

\begin{grammarlens}[title={What you've built}]
One more word for telling, explaining and describing at B1.
\end{grammarlens}

\subsection*{Your turn}
\begin{itemize}
\raggedright
  \item \emph{Say it:} \emph{das Darlehen}
  \item \emph{Say it:} \emph{das Darlehen}, once more
  \item \emph{Say it:} a sentence of your own with \emph{das Darlehen}
  \item \emph{Recall:} say the German for the amount, then the German for the bank loan, then say \emph{das Darlehen} again
\end{itemize}

\begin{tcolorbox}[breakable,colback=teal!4,colframe=teal!35!black,title={Before you move on}]
What does das Darlehen mean? (\textquotedblleft{}The loan\textquotedblright{}.) Say it once more.
\end{tcolorbox}
\section[die Schulden]{die Schulden — the debts}

\label{lesson:GE-C388-schulden}



\begin{quote}
Before the new one: say the German for the bank loan, then the German for the loan.
\end{quote}

\subsection*{You'll want to know: die Schulden}
\textbf{die Schulden} — \textquotedblleft{}the debts\textquotedblright{}.

\begin{quote}
Nach dem Autokauf hatte er hohe Schulden.
\end{quote}

\begin{quote}
\emph{After buying the car he had large debts.}
\end{quote}

A plural noun: \emph{Schulden haben} — to be in debt. The singular \emph{die Schuld} means guilt or blame.

\begin{cousinweb}
\emph{Schulden} is the plural of \emph{Schuld} \textquotedblleft{}debt, guilt\textquotedblright{}, inherited from Old High German \emph{sculd} and Proto-Germanic \emph{skuldiz} (EtymDB): what you owe.
\end{cousinweb}

\begin{grammarlens}[title={What you've built}]
One more word for telling, explaining and describing at B1.
\end{grammarlens}

\subsection*{Your turn}
\begin{itemize}
\raggedright
  \item \emph{Say it:} \emph{die Schulden}
  \item \emph{Say it:} \emph{die Schulden}, once more
  \item \emph{Say it:} a sentence of your own with \emph{die Schulden}
  \item \emph{Recall:} say the German for the bank loan, then the German for the loan, then say \emph{die Schulden} again
\end{itemize}

\begin{tcolorbox}[breakable,colback=teal!4,colframe=teal!35!black,title={Before you move on}]
What does die Schulden mean? (\textquotedblleft{}The debts\textquotedblright{}.) Say it once more.
\end{tcolorbox}
\section[zurückzahlen]{zurückzahlen — to pay back}

\label{lesson:GE-C388-zurueckzahlen}



\begin{quote}
Before the new one: say the German for the loan, then the German for the debts.
\end{quote}

\subsection*{You'll want to know: zurückzahlen}
\textbf{zurückzahlen} — \textquotedblleft{}to pay back\textquotedblright{}.

\begin{quote}
Ich zahle dir das Geld nächste Woche zurück.
\end{quote}

\begin{quote}
\emph{I will pay you the money back next week.}
\end{quote}

Separable: \emph{ich zahle zurück}, \emph{ich habe zurückgezahlt}.

\begin{cousinweb}
\emph{zurückzahlen} is built from \emph{zurück} \textquotedblleft{}back\textquotedblright{} and \emph{zahlen} \textquotedblleft{}to pay\textquotedblright{}, from \emph{Zahl} \textquotedblleft{}number\textquotedblright{}, Old High German \emph{zala}, Proto-Germanic \emph{talō} (EtymDB): you count the money back.
\end{cousinweb}

\begin{grammarlens}[title={What you've built}]
One more word for telling, explaining and describing at B1.
\end{grammarlens}

\subsection*{Your turn}
\begin{itemize}
\raggedright
  \item \emph{Say it:} \emph{zurückzahlen}
  \item \emph{Say it:} \emph{zurückzahlen}, once more
  \item \emph{Say it:} a sentence of your own with \emph{zurückzahlen}
  \item \emph{Recall:} say the German for the loan, then the German for the debts, then say \emph{zurückzahlen} again
\end{itemize}

\begin{tcolorbox}[breakable,colback=teal!4,colframe=teal!35!black,title={Before you move on}]
What does zurückzahlen mean? (\textquotedblleft{}To pay back\textquotedblright{}.) Say it once more.
\end{tcolorbox}
\section[borgen]{borgen — to borrow (sich borgen); to lend}

\label{lesson:GE-C388-borgen}



\begin{quote}
Before the new one: say the German for the debts, then the German for to pay back.
\end{quote}

\subsection*{You'll want to know: borgen}
\textbf{borgen} — \textquotedblleft{}to borrow (sich borgen); to lend\textquotedblright{}.

\begin{quote}
Ich habe mir das Fahrrad von meinem Nachbarn geborgt.
\end{quote}

\begin{quote}
\emph{I borrowed the bike from my neighbour.}
\end{quote}

\emph{Sich etwas borgen} — to borrow; \emph{jemandem etwas borgen} — to lend. In Austria people also say \emph{ausborgen}.

\begin{cousinweb}
\emph{borgen} is inherited from Old High German \emph{borgēn} and Proto-Germanic \emph{burgōną} (EtymDB), usually explained as first \textquotedblleft{}to look after, to give surety\textquotedblright{}, then to lend and borrow.
\end{cousinweb}

\begin{grammarlens}[title={What you've built}]
One more word for telling, explaining and describing at B1.
\end{grammarlens}

\subsection*{Your turn}
\begin{itemize}
\raggedright
  \item \emph{Say it:} \emph{borgen}
  \item \emph{Say it:} \emph{borgen}, once more
  \item \emph{Say it:} a sentence of your own with \emph{borgen}
  \item \emph{Recall:} say the German for the debts, then the German for to pay back, then say \emph{borgen} again
\end{itemize}

\begin{tcolorbox}[breakable,colback=teal!4,colframe=teal!35!black,title={Before you move on}]
What does borgen mean? (\textquotedblleft{}To borrow (sich borgen); to lend\textquotedblright{}.) Say it once more.
\end{tcolorbox}
